\documentclass[a4paper,12pt]{article}
\usepackage[utf8]{inputenc}
\usepackage[T1]{fontenc}
\usepackage[ngerman]{babel}
\usepackage{geometry}
\usepackage{amsmath}
\usepackage{amssymb}
\usepackage{hyperref}
\usepackage{booktabs}
\usepackage{array}
\usepackage{xcolor}
\geometry{margin=2.5cm}

\hypersetup{
    colorlinks=true,
    linkcolor=blue!60!black,
    urlcolor=blue!60!black,
    citecolor=blue!60!black
}

\begin{document}

% ================================================================
\section*{A Neural Network Multigrid Solver for the Navier--Stokes Equations}

\textbf{Autoren:} Nils Margenberg, Dirk Hartmann, Christian Lessig, Thomas Richter \\
\textbf{Zeitschrift:} Journal of Computational Physics, Band 460, 2022 \\
\textbf{DOI:} \texttt{10.1016/j.jcp.2022.110983}

% ================================================================
\subsection*{1\quad Problemstellung und Motivation}

Die direkte numerische Simulation der Navier-Stokes-Gleichungen zählt zu den
rechenintensivsten Aufgaben in der Ingenieur- und Naturwissenschaft. Klassische
Mehrgitterverfahren erzielen zwar eine optimal-lineare Gesamtkomplexität von
$\mathcal{O}(N)$, jedoch verursachen die mehrfach erforderlichen Glättungsschritte
auf feinen Gitterhierarchieebenen einen erheblichen Rechenaufwand pro Zeitschritt.
Margenberg et al.\ begegnen diesem Engpass, indem sie die teuersten Berechnungen auf
den feinsten Gitterebenen durch ein kompaktes neuronales Netz ersetzen.
Diese Hybridstrategie zielt darauf ab, den Effizienzgewinn von Lernmethoden
zu erschließen, ohne auf die mathematische Konsistenz des Mehrgitterrahmens
zu verzichten.

Die Fragestellung ist damit direkt verwandt mit der Motivation der vorliegenden
Masterarbeit: Auch dort stellt der konventionelle numerische Löser~-- in diesem
Fall LaMEM (Lithosphere and Mantle Evolution Model) für inkompressible
Stokes-Strömungen~-- einen Berechnungsengpass dar, der systematische
Parameterstudien für Mehr-Kristall-Sedimentationssysteme praktisch unmöglich
macht. Beide Arbeiten verfolgen das Ziel, durch datengetriebene Lernmethoden
eine Beschleunigung um mehrere Größenordnungen zu erzielen, ohne dabei die
physikalische Korrektheit der vorhergesagten Strömungsfelder wesentlich einzubüßen.

% ================================================================
\subsection*{2\quad Methodik: Der DNN-MG-Ansatz}

\paragraph{Grundprinzip.}
Der \emph{Deep Neural Network Multigrid Solver} (DNN-MG) erweitert ein klassisches
geometrisches Mehrgitterverfahren um ein neuronales Korrekturnetz. Das
Mehrgitterverfahren löst das lineare Gleichungssystem auf einer Hierarchie von
Gittern $\mathcal{T}_0 \subset \mathcal{T}_1 \subset \cdots \subset \mathcal{T}_L$,
wobei hochfrequente Fehleranteile auf feinen Ebenen geglättet und
niederfrequente Anteile auf groben Gittern korrigiert werden.
Auf den feinsten Ebenen $L+1$ (bzw.\ $L+2$) ersetzt ein neuronales Netz
den numerisch aufwändigen Mehrgitterschritt. Es sagt eine additive Korrektur
des prolongierten Geschwindigkeitsfeldes vorher:
\begin{equation}
    \tilde{v}^{L+1}_n = P\!\left(v^L_n\right) + d^{L+1}_n,
    \label{eq:dnnmg}
\end{equation}
wobei $P$ der Prolongationsoperator, $v^L_n$ die Mehrgitterlösung auf Ebene
$L$ und $d^{L+1}_n$ die netzpredizierte Korrektur zum Zeitschritt $n$ bezeichnet.
Die Korrektur $d^{L+1}_n$ fließt als Randbedingung in den nachfolgenden
Zeitschritt zurück, sodass das neuronale Netz zeitliche Kohärenz aktiv
mitmodelliert.

\paragraph{Netzwerkarchitektur.}
Das Korrekturmodell ist bewusst kompakt gehalten und operiert \emph{patch-basiert}:
Es verarbeitet ausschließlich kleine, lokale Ausschnitte (\emph{Patches}) des
Rechengebiets unabhängig voneinander, was die Parameterzahl niedrig hält und
gleichzeitig eine Verallgemeinerung auf veränderte Gebietsgeometrien erleichtert.
Als Eingangsgrößen empfängt das Netz:
\begin{itemize}
    \item die prolongierte Grobgitterlösung $P(v^L_n)$,
    \item nichtlineare Residuen der Navier-Stokes-Gleichungen auf Ebene $L+1$,
    \item lokale Péclet-Zahlen als Strömungsregimeindikatoren, sowie
    \item geometrische Informationen (Zellkantenlängen, Seitenverhältnisse, Winkel).
\end{itemize}
Als Kernbaustein wird eine \emph{Gated Recurrent Unit} (GRU) eingesetzt,
die temporale Abhängigkeiten zwischen aufeinanderfolgenden Zeitschritten
modelliert und so physikalisch kohärente Korrekturen über die Zeit hinweg
erzeugt. Das Gesamtmodell umfasst lediglich etwa $8\,600$ trainierbare Parameter,
was Training auf einem einzigen Mehrgitter-Datensatz ermöglicht.

\paragraph{Training und Evaluierung.}
DNN-MG wird am Benchmark-Problem der laminaren Kanalströmung um einen
elliptischen Zylinder bei Reynoldszahlen $Re = 100$ bis $200$ getestet.
Die Methode wird sowohl hinsichtlich Strömungsgenauigkeit (Lift- und
Widerstandsfunktionale, Wirbelablösefrequenz) als auch Recheneffizienz
(Zeit pro Zeitschritt) mit der reinen Mehrgitterlösung auf Ebene $L$ und
der vollen Lösung auf Ebene $L+1$ verglichen. Dabei erzielt DNN-MG:
\begin{itemize}
    \item eine Genauigkeit, die der vollständigen Mehrgitterlösung auf
          Ebene $L+1$ nahekommt,
    \item bei einem Rechenaufwand, der gegenüber der vollen $L+1$-Lösung
          um rund den Faktor~2 reduziert ist,
    \item und eine signifikante Verbesserung der Lift- und
          Widerstandsfunktionale gegenüber der reinen Grobgitterlösung auf Ebene $L$.
\end{itemize}
Besonders bemerkenswert ist, dass die Generalisierungsexperimente zeigen:
Ein auf \emph{einer} Hindernisgeometrie trainiertes Netz erzielt akzeptable
Ergebnisse für Kanalströmungen mit keinem oder mit zwei Hindernissen,
ohne auf diesen Geometrien nachtrainiert zu werden.

% ================================================================
\subsection*{3\quad Bezug zur Masterarbeit}

Die Arbeit von Margenberg et al.\ ist aus mehreren, komplementären
Perspektiven direkt relevant für die vorliegende Masterarbeit.
Die nachfolgenden Abschnitte vertiefen diese Verbindungen
systematisch unter Bezug auf Kapitel~3--5 der Masterarbeit.

% ---
\paragraph{3.1\quad Residuelles Lernen als gemeinsames Strukturprinzip.}

DNN-MG lernt in Gleichung~\eqref{eq:dnnmg} explizit eine
\emph{Korrektur} $d^{L+1}_n$ zur prolongierten Grobgitterlösung,
anstatt das vollständige Geschwindigkeitsfeld $\tilde{v}^{L+1}_n$
aus dem Nichts zu rekonstruieren.
Dieses residuelle Lernprinzip der Form
\begin{equation}
    v_\text{gesamt} = v_\text{Baseline} + \Delta v_\text{gelernt}
\end{equation}
ist strukturell analog zu Residual-Learning-Strategien, die auch in der
vorliegenden Masterarbeit als Erweiterungsoption diskutiert werden.
Dort dient die analytische Stokes-Lösung für einen einzelnen Zylinder
als physikalisch motivierte Basislösung, während das neuronale Netz
ausschließlich die verbleibende, durch Mehr-Körper-Wechselwirkungen
verursachte Abweichung erlernt. Beiden Ansätzen liegt die Erkenntnis
zugrunde, dass das Erlernen von Restfeldern numerisch stabiler ist
und eine stärkere physikalische Regularisierung bewirkt als die direkte
Schätzung des vollständigen Lösungsfeldes -- insbesondere wenn die Baseline
bereits einen Großteil der Lösungsstruktur korrekt beschreibt.

In der vorliegenden Arbeit wird als primäre Lernzielgröße stattdessen
die \emph{Stromfunktion} $\psi$ eingesetzt, die Inkompressibilität
analytisch durch $\nabla \cdot u = 0$ erzwingt. Dies stellt eine
verwandte, aber mathematisch rigorosere Form der physikalischen
Einbettung dar: Statt eine Korrektur zur numerischen Basislösung zu lernen,
wird die Lösungsdarstellung selbst so gewählt, dass ein zentrales physikalisches
Gesetz -- die Divergenzfreiheit des Geschwindigkeitsfeldes -- per Konstruktion
erfüllt ist.

% ---
\paragraph{3.2\quad Mehrskalenstruktur, U-Net-Architektur und Encoder-Decoder-Prinzip.}

Die hierarchische Gitterstruktur des Mehrgitterverfahrens~-- sukzessive
Vergröberung der Auflösung mit Fehlerglättung auf jeder Ebene, gefolgt von
Prolongation zurück auf das feine Gitter~-- ist strukturell äquivalent zur
Encoder-Decoder-Architektur des U-Nets, das in der vorliegenden Masterarbeit
als eine der beiden zentralen Surrogatarchitekturen eingesetzt wird.

Konkret lassen sich die Analogien wie folgt präzisieren:
\begin{center}
\begin{tabular}{lll}
\toprule
\textbf{Mehrgitterverfahren} & \textbf{U-Net} & \textbf{Funktion} \\
\midrule
Restriktion auf Grobgitter & Encoder-Pfad (Strided Conv.) & Auflösungsreduktion \\
Glättungsoperator auf jeder Ebene & Konvolutionsblöcke & Fehlerglättung \\
Prolongation auf feines Gitter & Decoder-Pfad (Transposed Conv.) & Auflösungserhöhung \\
Fehlerkorrektur von grob nach fein & Skip Connections & Informationsübertragung \\
\bottomrule
\end{tabular}
\end{center}
Beide Verfahren kombinieren Informationen aus verschiedenen räumlichen
Auflösungen: Das Mehrgitterverfahren durch Restriktion und Prolongation,
das U-Net durch Skip-Connections zwischen Encoder- und Decoderkanälen.
DNN-MG liefert damit eine numerisch fundierte Motivation für die Wahl
einer U-Net-basierten Architektur zur Vorhersage von Stokes-Strömungsfeldern
in der Masterarbeit: Die Mehrskalenstruktur ist kein rein empirisch motiviertes
Design, sondern entspricht einem tiefen strukturellen Zusammenhang zwischen
dem Charakter der zugrunde liegenden Differentialgleichung und der
verarbeitenden Architektur.

% ---
\paragraph{3.3\quad Generalisierung auf veränderte Kristall- und Hindernisanzahlen.}

Ein zentrales und für die Masterarbeit hochrelevantes Ergebnis von
Margenberg et al.\ ist die Fähigkeit des DNN-MG, auf Geometriekonfigurationen
zu generalisieren, die nicht im Training enthalten waren: Ein auf einem
einzigen Hindernis trainiertes Netz erzielt akzeptable Ergebnisse
für Strömungen mit null oder zwei Hindernissen.

Dies steht in direkter Analogie zur Kernfrage der vorliegenden
Masterarbeit, nämlich ob U-Net und Fourier-Neuronaler-Operator (FNO),
trainiert auf Konfigurationen mit bis zu $N_\text{max}$ Kristallen,
auf Konfigurationen mit einer variablen Anzahl zwischen $1$ und
$N_\text{max} + m$ Kristallen generalisieren können.

Margenberg et al.\ argumentieren, dass die \emph{patch-basierte},
lokal operierende Verarbeitung entscheidend für diese
Geometrieagnostizität ist: Da das Netz keine globale Gitterstruktur
explizit einlernt und stattdessen lokal reproduzierbare Muster
extrahiert, bleibt es gegenüber globalen Geometrieänderungen robust.

Für die vorliegende Masterarbeit ergibt sich daraus eine wichtige
Übertragungshypothese: Die geometrieagnostische Eingaberepräsentation
durch binäre Kristallmasken und vorzeichenbehaftete Abstandsfelder
(Signed Distance Fields) könnte eine vergleichbare Flexibilität
gegenüber variierenden Kristallzahlen ermöglichen~-- sofern das
Netz lokale Strömungsmuster in der Nachbarschaft einzelner Kristalle
zuverlässig extrahiert und diese zur Gesamtlösung überlagert.
Diese Hypothese ist eine der zentralen offenen Fragen, die in der
Masterarbeit empirisch untersucht wird.

% ---
\paragraph{3.4\quad Vergleich von lokalem und globalem Operatorlernen:
U-Net versus Fourier-Neuronaler-Operator.}

DNN-MG ist ein rein lokal operierendes Modell (patch-basiert),
das keine globalen spektralen Darstellungen des Strömungsfeldes verwendet.
Die vorliegende Masterarbeit ergänzt diese Perspektive durch den
systematischen Vergleich mit dem Fourier-Neuronalen-Operator (FNO),
der in einem einzigen Schicht die gesamte räumliche Domäne über
Fourier-Koeffizienten erfasst:
\begin{equation}
    v^{l+1}(x) = \sigma\!\left(
        W^l v^l(x) + \mathcal{F}^{-1}\!\left[R^l \cdot \mathcal{F}(v^l)\right](x)
    \right),
\end{equation}
wobei $\mathcal{F}$ die Fourier-Transformation, $R^l$ ein lernbarer
Gewichtstensor im Frequenzraum und $W^l$ eine punktweise lineare
Transformation darstellt.

Für elliptische Operatoren wie die Stokes-Gleichungen, bei denen
lokale Störungen durch das gesamte Strömungsfeld propagieren, ist
ein globaler Rezeptiv-Bereich pro Schicht strukturell vorteilhaft.
Das U-Net erreicht diesen globalen Kontext erst durch die Komposition
mehrerer Encoder-Stufen (hierarchisches Vergröbern), während der FNO
ihn direkt in jeder Schicht über den Spektralbereich zugänglich macht.
Margenberg et al.\ liefern damit ein theoretisches Argument,
das die lokale Mehrgittermethodik mit dem U-Net verbindet;
die Masterarbeit erweitert diesen Vergleichsrahmen, indem sie
beide Paradigmen~-- lokale Konvolution (U-Net) und globale
Spektraloperationen (FNO)~-- direkt auf demselben Datensatz
und mit identischem Trainingsprotokoll gegenüberstellt.

% ---
\paragraph{3.5\quad Physikalische Einbettung durch Residuen versus Lernzieldarstellung.}

DNN-MG integriert physikalische Information explizit als Eingabe: Das nichtlineare
Residuum der Navier-Stokes-Gleichungen auf Ebene $L+1$ dient als primäre
Eingabegröße. Dadurch wird das Netz auf physikalisch konsistente Korrekturen
konditioniert: Es muss nicht lernen, wann und wo Korrekturen physikalisch
notwendig sind, da diese Information direkt als Eingabesignal bereitsteht.

Die Masterarbeit verfolgt eine komplementäre Strategie physikalischer
Einbettung: Physikalisches Vorwissen wird nicht durch zusätzliche
Eingabekanäle, sondern durch die Wahl der \emph{Ausgabegröße}
injiziert. Durch die Vorhersage der Stromfunktion $\psi$ anstelle
der Geschwindigkeitskomponenten $u_x, u_z$ direkt wird die
Inkompressibilitätsbedingung
\begin{equation}
    \nabla \cdot u = 0 \quad \Longleftrightarrow \quad
    u_x = \frac{\partial \psi}{\partial z},\quad u_z = -\frac{\partial \psi}{\partial x}
\end{equation}
analytisch erzwungen, ohne dass sie als Verlustterm approximiert
werden müsste. Dies reduziert die effektiven Freiheitsgrade des
Lernproblems und verbessert die numerische Konditionierung der
Optimierung~-- beides Eigenschaften, die auch im DNN-MG-Kontext
durch die explizite Residuen-Konditionierung angestrebt werden.

Beide Strategien illustrieren dasselbe übergeordnete Prinzip:
Physikalisches Vorwissen als \emph{induktiver Bias} in das
Lernproblem einzubringen erhöht Dateneffizienz und physikalische
Konsistenz gleichzeitig. Sie unterscheiden sich lediglich darin,
\emph{wo} dieses Wissen injiziert wird: im Eingaberaum
(DNN-MG, über Residuen) oder im Ausgaberaum (Masterarbeit,
über Stromfunktion).

% ---
\paragraph{3.6\quad Wesentliche Unterschiede und Abgrenzung.}

Trotz der konzeptuellen Nähe bestehen bedeutsame Unterschiede
zwischen DNN-MG und dem in der Masterarbeit verfolgten Ansatz:

\begin{itemize}
    \item \textbf{Strömungsregime:} DNN-MG ist für instationäre
    Navier-Stokes-Strömungen bei moderaten Reynoldszahlen
    ($Re \approx 100$--$200$) konzipiert, bei denen Trägheitskräfte
    eine wesentliche Rolle spielen. Die vorliegende Masterarbeit
    behandelt ausschließlich stationäre Stokes-Strömungen
    ($Re \ll 1$) im magmatischen Kontext, bei denen Viskositätskräfte
    vollständig dominieren und das Strömungsfeld linear von den
    Randbedingungen und Kristallkonfigurationen abhängt.

    \item \textbf{Surrogat versus Hybridlöser:} DNN-MG ist ein
    Hybridlöser: Das neuronale Netz kann nicht eigenständig ohne
    einen begleitenden Mehrgitterlöser ausgeführt werden.
    Im Gegensatz dazu zielt die Masterarbeit auf ein
    vollständiges Surrogat ab, das~-- nach einmaligem Training~--
    für beliebige Kristallkonfigurationen eingesetzt wird,
    ohne dass eine numerische Simulation erforderlich ist.
    Dieses vollständige Surrogat-Paradigma ist besonders relevant
    für systematische Parameterstudien, in denen tausende
    Konfigurationen bewertet werden müssen.

    \item \textbf{Topologische Vielfalt:} DNN-MG wird auf
    Konfigurationen mit null, einem oder zwei Hindernissen
    untersucht~-- eine verhältnismäßig geringe topologische
    Variabilität. Die vorliegende Masterarbeit untersucht
    Konfigurationen mit bis zu $N_\text{max}$ zufällig
    positionierten und unterschiedlich großen Kristallen,
    was eine wesentlich komplexere topologische Vielfalt
    und ein fundamentaleres Generalisierungsproblem darstellt.

    \item \textbf{Architekturraum:} DNN-MG verwendet eine
    GRU-basierte Architektur für zeitliche Korrekturen.
    Die Masterarbeit vergleicht demgegenüber U-Net
    und FNO als zwei konzeptuell verschiedene Architekturen
    für stationäre Feldvorhersage -- einen ortsbasierten
    hierarchischen Ansatz gegen einen globalen Spektralansatz.
\end{itemize}

% ================================================================
\subsection*{4\quad Fazit und Einordnung}

Margenberg et al.\ demonstrieren überzeugend, dass die gezielte Kombination
klassischer numerischer Struktur mit kompakten neuronalen Netzen sowohl
wesentliche Effizienzgewinne als auch bemerkenswerte Generalisierungsfähigkeiten
ermöglicht. Insbesondere zeigen sie, dass:
\begin{enumerate}
    \item residuelles Lernen gegenüber vollständiger Feldvorhersage physikalisch
          robustere und effizientere Modelle liefert,
    \item die hierarchische Mehrskalenstruktur von Mehrgitterverfahren einen
          strukturellen Anknüpfungspunkt für U-Net-Architekturen bildet, und
    \item lokale, geometrieagnostische Verarbeitung entscheidend für die
          Generalisierung auf ungesehene Geometriekonfigurationen ist.
\end{enumerate}

Diese Erkenntnisse stärken direkt die methodischen Grundlagen der vorliegenden
Masterarbeit und motivieren insbesondere:
\begin{itemize}
    \item die Wahl der Stromfunktion~$\psi$ als physikalisch strukturierter
          Lernzielgröße (in Analogie zur residuellen Physikeinbettung
          in DNN-MG),
    \item die geometrieagnostische Eingabekodierung über Masken und
          Abstandsfelder (in Analogie zur patch-basierten lokalen
          Verarbeitung), sowie
    \item den systematischen Vergleich von U-Net und FNO als
          multiskaligen versus spektralen Architekturansätzen
          für Mehr-Kristall-Stokes-Strömungsfelder.
\end{itemize}

Die vorliegende Masterarbeit geht dabei über den von Margenberg et al.\
untersuchten Rahmen hinaus: Sie betrachtet ein vollständiges Surrogat
für stationäre Stokes-Strömungen, eine wesentlich größere topologische
Konfigurationsvielfalt, und führt einen direkten Vergleich zweier
fundamental verschiedener Lernparadigmen durch~-- ein Beitrag,
der in der bisherigen Literatur zu neuronalen Strömungssurrogaten
noch nicht systematisch adressiert wurde.

\end{document}